\documentclass[11pt,a4paper]{exam}
\usepackage[utf8]{inputenc}
\usepackage{amsmath, amsthm, amssymb, amsfonts} %ams packages
\usepackage{graphicx, color} %graphics
\usepackage{enumitem} %personalized lists 
\usepackage[top=3.5cm, bottom=2cm, left=2cm, right=2cm, headsep = 3pt, footskip=2pt]{geometry} %pagelayout
\usepackage{multirow,multicol} %to use multirow on tables
\usepackage{tabularx}
\usepackage{array} %bmatrix and similar stuff
% \usepackage{hyperref}
\usepackage[slovene]{babel}

\usepackage{multicol}


%%%%%%%%%%%%%%%%%%%%%%%%%%%%%%%%%%%%%%%%%%%%%%%%%%%%%%%%%%%%%%%%%%%%%
%personalized commands to make latex look better (optional)
%%%%%%%%%%%%%%%%%%%%%%%%%%%%%%%%%%%%%%%%%%%%%%%%%%%%%%%%%%%%%%%%%%%%%%%%%%%
\renewcommand{\leq}{\leqslant} %menores y mayores iguales decentes
\renewcommand{\geq}{\geqslant}
\renewcommand{\subset}{\subseteq} 
\renewcommand{\supset}{\supseteq} 
\renewcommand{\subsetneq}{\varsubsetneq}
\renewcommand{\{}{\lbrace}
\renewcommand{\}}{\rbrace}
\newcommand{\sm}{\setminus} %setminus corto

%\renewcommand{\bar}{\overline}
%\renewcommand{\hat}{\widehat}

%%%%%%%%%%%%%%%%%%%%%%%%%%%%%%%%%%%%%%%%%%%%%%%%%%%%%%%%%%%%%%%%%%%%%}

%%%%%%%%%%%%%%%%%%%%%%%%%%%%%%%%%%%%%%%%%%%%%%%%%%%%%%%%%%%%%%%%%%%%%%%%%%%
%exam class commands
%%%%%%%%%%%%%%%%%%%%%%%%%%%%%%%%%%%%%%%%%%%%%%%%%%%%%%%%%%%%%%%%%%%%%%%%%%%
\pagestyle{head}
% \chead{ Matematična analiza 2024/2025 \\ Univerza v Ljublani, Pedagoška fakulteta \\ Vaje 1 \\17. februar 2025} %self-explanatory
% \runningheader{}{ Abstraktna algebra 2024/2025 \\ Univerza v Ljublani, Pedagoška fakulteta \\2. Kolokvij \\20. januar 2025}{} %self-explanatory
% \runningheader{}{}{}
% \footer{}{}{}
% \firstpagefooter{MATH 1190 3.0 M  }{Page \thepage\ of \numpages}{Copyright 2020 \textcopyright Antonio Montero}
% \firstpagefooter{}{}{}
% \runningfooter{MATH 1190 3.0 M }{Page \thepage\ of \numpages}{Copyright 2020 \textcopyright Antonio Montero}
% \firstpagefootrule 
% \runningfootrule

%%Solutions
\unframedsolutions
% \printanswers
\SolutionEmphasis{\itshape}


%Personalized commands
\pointpoints{točka}{točk}
\bonuspointpoints{točka}{dodatnih točk}
% \gradetable[h][questions]
\addpoints
\hqword{Naloga}
% \vpgword{text}
\hpword{Točke}
% \hsword{}
% \hbpword{}
\htword{Skupaj}

% \pointsinrightmargin
% \extrawidth{-2cm}
\marginpointname{ \points}


\chqword{Naloga}
\chpgword{Stran}
\chpword{točke}
% \chbpword{Puntos extra}
% \chsword{Puntos obtenidos
\chtword{Skupaj}
%%%%%%%%%%%%%%%%%%%%%%%%%%%%%%%%%%%%%%%%%%%%%%%%%%%%%%%%%%%%%%%%%%%%%

% \graphicspath{{../img/}} %Uso: \graphicspath{{RelativePath}}

%%%%%%%%%%%%%%%%%%%%%%%%%%%%%%%%%%%%%%%%%%%%%%%%%%%%%%%%%%%%%%%%%%%%%
%Personalized commands

%Personalized commands

\usepackage{tabularx}
\newcolumntype{C}{>{\centering\arraybackslash}X}%
\newcolumntype{R}{>{\raggedleft\arraybackslash}X}%
\newcolumntype{L}{>{\raggedright\arraybackslash}X}%

\newcolumntype{\cc}[1]{>{\hsize=#1\hsize\raggedright\arraybackslash}X}%
\newcolumntype{\rr}[1]{>{\hsize=#1\hsize\raggedleft\arraybackslash}X}%
\newcolumntype{\ll}[1]{>{\hsize=#1\hsize\centering\arraybackslash}X}%

\newcommand{\TT}{\mathbf{T}}
\newcommand{\FF}{\mathbf{F}}

\DeclareMathOperator{\sh}{sh}
\DeclareMathOperator{\ch}{ch}
\let\th\relax
\DeclareMathOperator{\th}{th}
\DeclareMathOperator{\arsh}{arsh}
\DeclareMathOperator{\arch}{arch}
\DeclareMathOperator{\arth}{arth}
\DeclareMathOperator{\arctg}{arctg}

%%%%%%%%%%%%%%%%%%%%%%%%%%%%%%%%%%%%%%%%%%%%%%%%%%%%%%%%%%%%%%%%%%%%%
		\newif\ifmarking
% 		\markingtrue %comment to hide marking suggestions

		\newenvironment{marking}
		{\hrulefill
		 
		\textbf{Marking suggestions:} 
		\ttfamily }
		{\hrulefill }

\allowdisplaybreaks

% \renewcommand{\thepartno}{\roman{partno}}
\renewcommand{\partlabel}{(\textit{\thepartno})}
%%%%%%%%%%%%%%%%%%%%%%%%%%%%%%%%%%%%%%%%%%%%%%%%%%%%%%%%%%%%%%%%%%%%%

\begin{document}
\chead{Matematična analiza 2024/2025 \\ Univerza v Ljubljani, Pedagoška fakulteta \\ 3. Izpit \\18. avgust 2025}
% \footer{}{}{}

% \vspace{1 cm}

\pagestyle{head}
	

	
\begin{questions}

\question
	Dana je funkcija $g(x) = \frac{x^{3}}{3}$.
	\begin{parts}
  \part Poiščite intervale naraščanja in padanja ter stacionarne točke funkcije $g$. Za stacionarne točke  določite, ali so lokalni/globalni minimum, maksimum ali nič od tega.
  \part Čimbolj natančno skicirajte graf funkcije $g$ s pomočjo prejšnje točke.
	\part Določite tangento in normalno na graf funkcije $g$ v točki $(-3,-9)$. 
	\part Poiščite vse tangente na graf funkcije $g$, ki so vzporedne premici $y=4x-2$, ter tudi vse tangente, ki sekajo $x$-os pod kotom $\frac{\pi}{4}$.
	\part Določite koeficient premice $y=kx-18$, da bo ta tangenta na graf funkcije $g$.
	\end{parts}

\question V enem obratu tovarne izdelujejo pločevinaste posode brez pokrova, ki so oblike valja z volumnom $1 \ \mathrm{dm}^3$, v drugem obratu te tovarne pa izdelujejo plastične pokrove za te posode. Cena $1 \text{ dm}^2$ pločevine je $2$ centa, cena $1 \ \mathrm{dm}^2$ plastike pa $1$ cent. Kakšne naj bodo dimenzije posode (s pokrovom), da bodo stroški materiala minimalni? (Če ne gre, določi dimenzije posode brez pokrova, da porabimo čim manj pločevine.)


\question Dana je funkcija 
\[
  f(x)=
  \begin{cases}
    3-x & \text{za }  0 \leq x < 1 \\
    -x^{2} & \text{za } 1 \leq x < 2 \\
    1 & \text{za } 2 \leq x \leq 3 \\
  \end{cases}
\]
	
\begin{parts}
		\part Naj bo $\mathcal{D}: 0 < \frac{1}{2} < 1 < \frac{5}{4} < \frac{8}{5} < 2 < \frac{7}{3} < \frac{11}{4} < 3 $ delitev intervala $[0,3]$. Zapišite zgornjo in spodnjo Darbouxjevo vsoto funkcije $f$ za dano delitev $\mathcal{D}$, ter ju izračunajte. 
    \part Za poljubno število $n \in \mathbb{N}$ naj bo $\mathcal{D}_{n}: 0 < \frac{1}{n} < \frac{2}{n} < \cdots < \frac{3(n-1)}{n} < 3$ delitev intervala $[0,3]$. Zapišite zgornjo Darbouxjevo vsoto $S(\mathcal{D}_{n})$ in spodnjo Darbouxjevo vsoto $s(\mathcal{D}_{n})$ funkcije $f$ za dano delitev $\mathcal{D}_{n}$.
\end{parts}

\question Izračunajte naslednje nedoločene integrale: 
  \begin{parts}
    \part  $\displaystyle \int x^{3} (\sin x) dx$;
    \part $\displaystyle \int \frac{(\ln x)^2}{x(1 + (\ln x)^2)} dx $;
    \part $\displaystyle \int \frac{dx}{x \left((\ln x)^2 -1 \right)} $.
  \end{parts}

\question Naj bo $\{f_n: [0,1] \to \mathbb{R}\}$ zaporedje funkcij, definirano z $f_n(x) = \frac{x^{n}}{1+x^{n}}$ za vsak $n \in \mathbb{N}$.
\begin{parts}
    \part Ugotovite, če  zaporedje $\{f_{n}\}$ konvergira po točkah in če konvergira, določite njegovo limitno funkcijo $f(x)$.
    \part Določite, če je zaporedje $\{f_{n}\}$ enakomerno konvergentno.
\end{parts}
\end{questions}

\end{document}
